\section{Conclusion}\label{sec:conclusion}
A certified neural Bellman method needs an approximation space compatible with the stopping problem, an admissible policy, and complete error inequalities tied to that same policy. The preference-adjustment economy illustrates why these requirements must be checked separately. Its inward upper-wealth boundary creates a strict discontinuity between continuation value and immediate liquidation. The accessibility-preserving architecture and signed facewise theorem remove the resulting trace budget without changing the economic model or restricting the upper comparison's control opportunity.

The new experiments establish distinct advances: complete multi-seed original-economy neural verification with exact trace control; six continuously verified state-space comparison policies; a state-global neural certificate in coupled inventory dynamics up to 128 dimensions; an established-arithmetic recalculation of every decisive historical price node; and fresh end-to-end generation of the full sharp price continuum. A quantitative policy-iteration recurrence states exactly which verified approximation defects imply convergence. A consumption-top-up bound supplies an economic interpretation with an explicit financing convention.

The objective remains an accurate and competitive neural policy/value solver for nonlinear dynamic economies. On the original full-horizon economy, the present neural residual certificates still do not reach $.01$, and the newly implemented classical candidates do not yet supply a matched-accuracy frontier. The structured inventory success and sharp time-control bounds do not replace those missing nonlinear accuracy results. They supply reproducible constructions, mathematically justified improvements, and separately identified comparison objects through which further numerical advances can be assessed without changing the question being solved.
